\section*{Conclusions}

Taxonomy must keep pace with the rapid development of digital biology. Although analytical methods have advanced considerably, species descriptions and manuscripts are still often prepared as conventional text, limiting their computational use and integration with other biodiversity data.

This study demonstrates that semantic species descriptions can be incorporated into a complete taxonomic revision and carried directly from their computable source into a publication-ready manuscript. The integrated use of Phenoscript, GitHub, and Overleaf provides a practical model for producing taxonomic treatments that remain readable to taxonomists while being available for computational reuse.

The revision also reveals previously undocumented diversity in the Eastern Arc Mountains by adding four new species of \textit{Amietina}. Their restricted distributions provide further evidence that separate mountain blocks harbour distinct dung beetle faunas and emphasize the importance of documenting the remaining rainforest fragments. Together, these methodological and taxonomic results show how structured digital approaches can support both the study of biodiversity and the preservation of the knowledge derived from it.
